\documentclass[12pt,letterpaper]{article}

\usepackage[spanish]{babel}
\usepackage[utf8]{inputenc}
\usepackage{amsmath, amssymb}
\usepackage{geometry}
\usepackage{graphicx}
\usepackage{fancyhdr}

\geometry{margin=2.5cm}
\setlength{\headheight}{52pt}

\pagestyle{fancy}
\fancyhf{}

\fancyhead[L]{
    \includegraphics[height=1.5cm]{ur_image.png}
}

\fancyhead[R}{
\begin{tabular}{r}
\textbf{Monitoria Álgebra Lineal - Ortogonalidad} \\
Gabriel Fernando Salcedo Zamora \\
\end{tabular}
}

\begin{document}

\section*{Taller: Ortogonalidad y proyecciones}

\begin{enumerate}

% Producto punto
\item Calcule el producto punto y determine si son ortogonales:
\[
(1,2,3), (2,-1,0)
\]

% Proyección
\item Encuentre la proyección de $\mathbf{u}=(2,1)$ sobre $\mathbf{v}=(1,1)$.

% Ortogonalización
\item Aplique Gram-Schmidt a:
\[
(1,1,0), (1,0,1)
\]

% Subespacio ortogonal
\item Encuentre una base de $W^\perp$ si:
\[
W = \text{span}\{(1,2,3)\}
\]

\end{enumerate}

\newpage

\section*{Solución}

\begin{enumerate}

\item Producto punto:
\[
1\cdot2 + 2(-1) + 3\cdot0 = 2-2+0=0.
\]

Son ortogonales.

\item 
\[
\text{proj}_{\mathbf{v}}\mathbf{u} =
\frac{\mathbf{u}\cdot \mathbf{v}}{\mathbf{v}\cdot \mathbf{v}}\mathbf{v}
\]

\[
= \frac{2+1}{1+1}(1,1) = \frac{3}{2}(1,1) = (3/2,3/2).
\]

\item 

\[
u_1 = (1,1,0)
\]

\[
u_2 = (1,0,1) - \frac{(1,0,1)\cdot u_1}{u_1\cdot u_1}u_1
\]

\[
= (1,0,1) - \frac{1}{2}(1,1,0)
= (1/2,-1/2,1)
\]

\item Buscamos:
\[
(x,y,z)\cdot(1,2,3)=0 \Rightarrow x+2y+3z=0.
\]

Base:
\[
\{(-2,1,0),(-3,0,1)\}.
\]

\end{enumerate}

\end{document}
